\subsection{Evaluation Details}
\label{sec:experiments_evaluationdetails}

MS-COCO is a standard evaluation set used in existing text-to-image works \cite{ramesh-dalle,xu-cvpr-2018,nichol-glide}. Hence, we perform majority of comparisons and ablation studies on MS-COCO. We also compare \name with recent works \cite{ramesh-dalle2, crowson2022vqgan, rombach-cvpr-2022} on our proposed benchmark \benchmarkname.

We evaluate both image quality and image-text alignment across several metrics; we use automated metrics as well as human evaluation. For automated metrics, we use FID score \cite{heusel2017gans} to evaluate the fidelity of the samples, and CLIP score \cite{hessel2021clipscore, radford2021learning} to measure the alignment between the generated image and input text prompt. FID score compares the distribution of Inception features between the generated samples and the reference images, with a lower FID score generally corresponding to higher fidelity. CLIP score measures the dot product between the encoded reference text and the generated image, a higher CLIP score would generally imply better alignment under the CLIP model. Since guidance weight is an important knob to control image quality and text alignment, we report most of our ablation results using trade-off curves between CLIP and FID scores across a range of guidance weights. Both metrics have limitations, for example FID is not fully aligned with perceptual quality \cite{parmar-cvpr-2022}, and CLIP is unable to perform counting well \cite{radford2021learning}. Due to these limitations, we also use human evaluation to assess for image quality, and caption similarity. 

To assess the image quality on MS-COCO, we run a 2-alternative forced choice (2AFC) evaluation where the rater is asked to choose between the model generation and the reference image based on the question: ``Which image is more photorealistic (looks more real)?''. Note that this set up is different from that of \benchmarkname, which does not have reference images. We report the percentage of times raters choose model's generations over reference images (a.k.a. the \emph{preference rate}). A 50\% preference rate implies the evaluator has equal preference between the model's samples vs. the reference.

For image-text alignment, human raters are shown a model-generated image (or a reference image) together with the input caption, and asked ``Does the caption accurately describe the above image?''. The raters are allowed to choose among three possible options : ``yes'', ``somewhat'', ``no'', each corresponding to scores of 100, 50, and 0, respectively. This was done independently for both model samples and the reference images. We report both values (e.g., the reference image may not be well aligned with the corresponding caption in certain cases). We also report evaluation results on a filtered set of examples, i.e. examples from the reference that do not contain any person in it. We filter based on a set of people-specific identifiers in text prompts.

For both types of human evaluations, we use 200 randomly chosen (image, caption) pairs from the MS-COCO test set, and ran the human survey in four batches each containing 50 images. We obtained total of 100 responses for each of 200 captions.
Before we aggregate the responses, we interleave 10 ``control'' questions and only count responses from raters who correctly rated at least 80\% of the control questions. After this filtering, we had at least 73 and 51 ratings per image for image quality and image-text alignment evaluations, respectively.

